
\documentclass{article}
\usepackage{colortbl}
\usepackage{makecell}
\usepackage{multirow}
\usepackage{supertabular}

\begin{document}

\newcounter{utterance}

\centering \large Interaction Transcript for game `cladder', experiment `full\_v1.5\_default', episode 9197 with refineonly{-}Qwen3.5{-}27B.
\vspace{24pt}

{ \footnotesize  \setcounter{utterance}{1}
\setlength{\tabcolsep}{0pt}
\begin{supertabular}{c@{$\;$}|p{.15\linewidth}@{}p{.15\linewidth}p{.15\linewidth}p{.15\linewidth}p{.15\linewidth}p{.15\linewidth}}
    \# & \multicolumn{2}{c}{Player} && \multicolumn{2}{c}{Game Master} \\
    \hline

    \theutterance \stepcounter{utterance}  
    & & & \multicolumn{4}{p{0.6\linewidth}}{
        \cellcolor[rgb]{0.9,0.9,0.9}{
            \makecell[{{p{\linewidth}}}]{
                \texttt{\tiny{[P1$\langle$GM]}}
                \texttt{You are an expert in causal inference. The following question is not a typical commonsense query, but rather a meticulously designed question created by a professor specializing in causal inference, intended to assess the students' mastery of the course content.} \\
\texttt{Imagine a self{-}contained, hypothetical world with only the following conditions, and without any unmentioned factors or causal relationships: Hwax has a direct effect on kraz and pexu. Pexu has a direct effect on rukz. Kraz has a direct effect on rukz. The overall probability of pexu is 82\%. The probability of not pexu and rukz is 13\%. The probability of pexu and rukz is 47\%. Is the chance of rukz smaller when observing pexu?} \\
\texttt{Start your answer with 'Yes' or 'No', followed by additional reasoning or evidence to support your explanation.} \\
            }
        }
    }
     \\ \\

    \theutterance \stepcounter{utterance}  
    & & \multicolumn{4}{p{0.6\linewidth}}{
        \cellcolor[rgb]{0.95,0.95,0.95}{
            \makecell[{{p{\linewidth}}}]{
                \texttt{\tiny{[P1$\leftrightarrow$P]}}
                \texttt{{[}DRAFT{]}} \\
\texttt{No} \\
\\ 
\texttt{To determine whether the chance of *rukz* is smaller when observing *pexu*, we must calculate and compare the conditional probability \$P(\textbackslash{}text\{rukz\} \textbackslash{}mid \textbackslash{}text\{pexu\})\$ with the marginal probability \$P(\textbackslash{}text\{rukz\})\$.} \\
\\ 
\texttt{First, let's derive the marginal probability of *rukz*, \$P(\textbackslash{}text\{rukz\})\$. We are given the joint probabilities for *rukz* occurring with and without *pexu*:} \\
\texttt{1.  \$P(\textbackslash{}text\{pexu\} \textbackslash{}cap \textbackslash{}text\{rukz\}) = 47\textbackslash{}\% = 0.47\$} \\
\texttt{2.  \$P(\textbackslash{}neg\textbackslash{}text\{pexu\} \textbackslash{}cap \textbackslash{}text\{rukz\}) = 13\textbackslash{}\% = 0.13\$} \\
\\ 
\texttt{Since *pexu* and \$\textbackslash{}neg\textbackslash{}text\{pexu\}\$ are mutually exclusive and exhaustive events, the total probability of *rukz* is the sum of these joint probabilities:} \\
\texttt{\$\$P(\textbackslash{}text\{rukz\}) = P(\textbackslash{}text\{pexu\} \textbackslash{}cap \textbackslash{}text\{rukz\}) + P(\textbackslash{}neg\textbackslash{}text\{pexu\} \textbackslash{}cap \textbackslash{}text\{rukz\})\$\$} \\
\texttt{\$\$P(\textbackslash{}text\{rukz\}) = 0.47 + 0.13 = 0.60\$\$} \\
\\ 
\texttt{Next, we calculate the conditional probability of *rukz* given *pexu*, \$P(\textbackslash{}text\{rukz\} \textbackslash{}mid \textbackslash{}text\{pexu\})\$. The formula for conditional probability is:} \\
\texttt{\$\$P(\textbackslash{}text\{rukz\} \textbackslash{}mid \textbackslash{}text\{pexu\}) = \textbackslash{}frac\{P(\textbackslash{}text\{pexu\} \textbackslash{}cap \textbackslash{}text\{rukz\})\}\{P(\textbackslash{}text\{pexu\})\}\$\$} \\
\\ 
\texttt{We are given that the overall probability of *pexu* is \$82\textbackslash{}\%\$ (\$0.82\$). Substituting the known values:} \\
\texttt{\$\$P(\textbackslash{}text\{rukz\} \textbackslash{}mid \textbackslash{}text\{pexu\}) = \textbackslash{}frac\{0.47\}\{0.82\} \textbackslash{}approx 0.573\$\$} \\
\\ 
\texttt{Finally, we compare the two probabilities:} \\
\texttt{*   \$P(\textbackslash{}text\{rukz\}) = 0.60\$} \\
\texttt{*   \$P(\textbackslash{}text\{rukz\} \textbackslash{}mid \textbackslash{}text\{pexu\}) \textbackslash{}approx 0.573\$} \\
\\ 
\texttt{Since \$0.573 < 0.60\$, the probability of *rukz* is indeed lower when *pexu* is observed compared to the overall average. This counter{-}intuitive result (where a direct cause *pexu* seems to lower the probability of its effect *rukz* in the marginal distribution) is a classic example of **Simpson's Paradox** or a specific instance} \\
            }
        }
    }
    & \\ \\

    \theutterance \stepcounter{utterance}  
    & \multicolumn{4}{p{0.6\linewidth}}{
        \cellcolor[rgb]{0.561,0.737,0.561}{
            \makecell[{{p{\linewidth}}}]{
                \texttt{\tiny{[P1$\rangle$GM]}}
                \texttt{No} \\
\\ 
\texttt{To determine whether the chance of *rukz* is smaller when observing *pexu*, we must calculate and compare the conditional probability \$P(\textbackslash{}text\{rukz\} \textbackslash{}mid \textbackslash{}text\{pexu\})\$ with the marginal probability \$P(\textbackslash{}text\{rukz\})\$.} \\
\\ 
\texttt{First, let's derive the marginal probability of *rukz*, \$P(\textbackslash{}text\{rukz\})\$. We are given the joint probabilities for *rukz* occurring with and without *pexu*:} \\
\texttt{1.  \$P(\textbackslash{}text\{pexu\} \textbackslash{}cap \textbackslash{}text\{rukz\}) = 47\textbackslash{}\% = 0.47\$} \\
\texttt{2.  \$P(\textbackslash{}neg\textbackslash{}text\{pexu\} \textbackslash{}cap \textbackslash{}text\{rukz\}) = 13\textbackslash{}\% = 0.13\$} \\
\\ 
\texttt{Since *pexu* and \$\textbackslash{}neg\textbackslash{}text\{pexu\}\$ are mutually exclusive and exhaustive events, the total probability of *rukz* is the sum of these joint probabilities:} \\
\texttt{\$\$P(\textbackslash{}text\{rukz\}) = P(\textbackslash{}text\{pexu\} \textbackslash{}cap \textbackslash{}text\{rukz\}) + P(\textbackslash{}neg\textbackslash{}text\{pexu\} \textbackslash{}cap \textbackslash{}text\{rukz\})\$\$} \\
\texttt{\$\$P(\textbackslash{}text\{rukz\}) = 0.47 + 0.13 = 0.60\$\$} \\
\\ 
\texttt{Next, we calculate the conditional probability of *rukz* given *pexu*, \$P(\textbackslash{}text\{rukz\} \textbackslash{}mid \textbackslash{}text\{pexu\})\$. The formula for conditional probability is:} \\
\texttt{\$\$P(\textbackslash{}text\{rukz\} \textbackslash{}mid \textbackslash{}text\{pexu\}) = \textbackslash{}frac\{P(\textbackslash{}text\{pexu\} \textbackslash{}cap \textbackslash{}text\{rukz\})\}\{P(\textbackslash{}text\{pexu\})\}\$\$} \\
\\ 
\texttt{We are given that the overall probability of *pexu* is \$82\textbackslash{}\%\$ (\$0.82\$). Substituting the known values:} \\
\texttt{\$\$P(\textbackslash{}text\{rukz\} \textbackslash{}mid \textbackslash{}text\{pexu\}) = \textbackslash{}frac\{0.47\}\{0.82\} \textbackslash{}approx 0.573\$\$} \\
\\ 
\texttt{Finally, we compare the two probabilities:} \\
\texttt{*   \$P(\textbackslash{}text\{rukz\}) = 0.60\$} \\
\texttt{*   \$P(\textbackslash{}text\{rukz\} \textbackslash{}mid \textbackslash{}text\{pexu\}) \textbackslash{}approx 0.573\$} \\
\\ 
\texttt{Since \$0.573 < 0.60\$, the probability of *rukz* is indeed lower when *pexu* is observed compared to the overall average. This result, where the conditional probability differs from the marginal probability despite a direct causal link, highlights the distinction between causal effects and associative probabilities in the presence of confounding variables (specifically, *hwax* acts as a common cause of both} \\
            }
        }
    }
    & & \\ \\

    \theutterance \stepcounter{utterance}  
    & & & \multicolumn{2}{p{0.3\linewidth}}{
        \cellcolor[rgb]{0.9,0.9,0.9}{
            \makecell[{{p{\linewidth}}}]{
                \texttt{\tiny{[GM$|$GM]}}
                \texttt{No} \\
            }
        }
    }
    & & \\ \\

    \theutterance \stepcounter{utterance}  
    & & & \multicolumn{2}{p{0.3\linewidth}}{
        \cellcolor[rgb]{0.9,0.9,0.9}{
            \makecell[{{p{\linewidth}}}]{
                \texttt{\tiny{[GM$|$GM]}}
                \texttt{yes} \\
            }
        }
    }
    & & \\ \\

    \theutterance \stepcounter{utterance}  
    & & & \multicolumn{2}{p{0.3\linewidth}}{
        \cellcolor[rgb]{0.9,0.9,0.9}{
            \makecell[{{p{\linewidth}}}]{
                \texttt{\tiny{[GM$|$GM]}}
                \texttt{game\_result = LOSE} \\
            }
        }
    }
    & & \\ \\

\end{supertabular}
}

\end{document}
